\documentclass[10pt]{article}
\usepackage[T1]{fontenc}
\usepackage[utf8]{inputenc}
\usepackage[french]{babel}
\usepackage{amsmath,amssymb,mathrsfs,mathpazo}
\usepackage{geometry}
\geometry{a4paper,margin=2cm}
\usepackage{graphicx,array,tabularx,tabularray,colortbl,multirow,multicol,booktabs,enumitem,pifont,ulem,fancybox,diagbox,cancel}
\usepackage{pstricks,pstricks-add,pst-plot,pst-tree,pst-node,pst-text,pst-func,pst-math,pst-bspline,pst-3dplot,pst-eucl,pst-solides3d}
\let\pscalculate\undefined
\usepackage{pst-bezier}
\usepackage{tikz,pgfplots,tkz-tab}
\usetikzlibrary{arrows,arrows.meta,patterns,calc,positioning,trees,decorations.markings,shapes.geometric,shapes.misc,shapes.arrows,intersections,decorations.pathmorphing,decorations.pathreplacing,backgrounds,shadows,babel}
\pgfplotsset{compat=1.18}
\usepgfplotslibrary{fillbetween}
\usepackage[np]{numprint}
\usepackage{listings,makecell,pas-tableur,twoopt}
\let\barre\undefined
\usepackage[pstricks]{tablvar}
\definecolor{darkgreen}{RGB}{0,100,0}
\providecommand{\red}{\color{red}}
\providecommand{\blue}{\color{blue}}
\providecommand{\green}{\color{green}}
\providecommand{\black}{\color{black}}
\providecommand{\white}{\color{white}}
\providecommand{\gray}{\color{gray}}
\providecommand{\cyan}{\color{cyan}}
\providecommand{\magenta}{\color{magenta}}
\providecommand{\yellow}{\color{yellow}}
\providecommand{\violet}{\color{violet}}
\providecommand{\orange}{\color{orange}}
\providecommand{\euro}{\texteuro}
\providecommand{\EUR}{\texteuro}
\providecommand{\ieme}{\textsuperscript{e}}
\providecommand{\degres}{\ensuremath{^\circ}}
\providecommand{\degree}{\ensuremath{^\circ}}
\providecommand{\vv}[1]{\overrightarrow{#1}}
\providecommand{\Coord}[1]{#1}
\providecommand{\decosix}{}
\providecommand{\starredbullet}{\ensuremath{\bullet}}
\pagestyle{empty}
\setlength{\parindent}{0pt}
\usepackage{adjustbox}\DeclareUnicodeCharacter{2236}{:}\DeclareUnicodeCharacter{2190}{\ensuremath{\leftarrow}}
\providecommand{\R}{\mathbb{R}}
\providecommand{\N}{\mathbb{N}}
\providecommand{\Z}{\mathbb{Z}}
\providecommand{\Q}{\mathbb{Q}}
\providecommand{\D}{\mathbb{D}}
\providecommand{\C}{\mathbb{C}}
\providecommand{\e}{{\mathrm{e}}}
\providecommand{\vect}[1]{\overrightarrow{#1}}
\providecommand{\vectt}[1]{\overrightarrow{\mathrm{#1}}}
\providecommand{\vvt}[1]{\overrightarrow{\mathrm{#1}}}
\providecommand{\np}[2][]{{#2~#1}}
\providecommand{\vv}[1]{{\overrightarrow{#1}}}
\renewcommand{\d}{{\mathrm{d}}}
\providecommand{\degres}{{^\circ}}
\providecommand{\degre}{{^\circ}}
\providecommand{\blue}{\color{blue}}
\providecommand{\red}{\color{red}}
\providecommand{\enskip}{\quad}
\providecommand{\llbracket}{\lbrack\!\lbrack}
\providecommand{\rrbracket}{\rbrack\!\rbrack}
\providecommand{\euro}{{\text{€}}}
\providecommand{\no}{{\text{n°}}}
\providecommand{\interval}[2]{{[#1;#2]}}
\providecommand{\intervFO}[2]{{[#1;#2[}}
\providecommand{\dotfill}{\ldots}

\providecommand{\ds}{\displaystyle}
\providecommand{\textding}[1]{#1}
\providecommand{\Oij}{\(\mathrm{O};\vec{\imath},\vec{\jmath}\)}
\providecommand{\Oijk}{\(\mathrm{O};\vec{\imath},\vec{\jmath},\vec{k}\)}
\providecommand{\Ouv}{\(\mathrm{O};\vec{u},\vec{v}\)}
\providecommand{\intervFF}[2]{\ensuremath{[#1;#2]}}
\renewcommand{\euro}{\ifmmode\text{\texteuro}\else\texteuro\fi}
\DeclareUnicodeCharacter{00B0}{\ifmmode^{\circ}\else\textdegree\fi}
\DeclareUnicodeCharacter{20AC}{\ifmmode\text{\texteuro}\else\texteuro\fi}
% Énoncés : Éducation nationale. Transcription et corrigés : APMEP, Denis Vergès et contributeurs. Découpage et classement : CoopMaths, Sébastien Lozano, Éric Elter et contributeurs. Adaptation de la présentation : Maths Explorer.
% CC BY-SA : https://coopmaths.fr/www/about/
% https://coopmaths.fr/alea/static/bac/2022/tex/bac_2022_09_sujet1_metropole_4_cor.tex
\begin{document}
\noindent\textit{APMEP / CoopMaths. CC BY-SA. Présentation adaptée par Maths Explorer.}\par\medskip
\medskip

Dans l'espace rapporté à un repère orthonormé \Oijk, on considère les points

\[\text{A}(-1~;~-1~;~3),\qquad \text{B}(1~;~1~;~2),\qquad  \text{C}(1~;~-1~;~7)\]

On considère également la droite $\Delta$ passant par les points D$(-1~;~6~;~8)$ et E$(11~;~- 9~;~2)$.

\medskip

\begin{enumerate}
\item 
	\begin{enumerate}
		\item 
		





Avec $\vect{\text{DE}}\begin{pmatrix}12\\-15\\-6\end{pmatrix}$, on peut prendre comme vecteur directeur de $\Delta$:
$\dfrac13\vect{\text{DE}}\begin{pmatrix}4\\-5\\-2\end{pmatrix}$.

On a  $M(x~;~y~;~z) \in (\text{DE}) \iff \text{il existe } t \in \R, \: \text{tel que } \: \vect{\text{D}M} = t \dfrac13\vect{\text{DE}}$, soit 

\[\left\{\begin{array}{l c r}
x + 1 &=& 4t\\
y - 6 &=&-5t\\
z - 8 &=& -2t
\end{array}\right. \quad\text{avec}\: t \in \R
\iff 
\left\{\begin{array}{l c l}
x &=& -1 + 4t\\
y &=& \phantom{-}6 - 5t \\
z &=& \phantom{-}8 - 2t
\end{array}\right. \quad\text{avec}\: t \in \R\]
		\item La droite $\Delta'$ est parallèle à $\Delta$ donc elle a les mêmes vecteurs directeurs que $\Delta$. De plus, elle passe par l'origine O du repère donc elle a pour représentation paramétrique:

$\left\{\begin{array}{l !{=} r}
x&0+4t\\
y&0-5t\\
z& 0-2t
\end{array}\right. \quad\text{avec}\: t \in \R$
soit
$\left\{\begin{array}{l !{=} r}
x&4t\\
y&-5t\\
z& -2t
\end{array}\right. \quad\text{avec}\: t \in \R$

		\item 
		F$(1,36~;~-1,7~;~-0,7) \in\Delta' \iff \left\{\begin{array}{l c l}
1,36&=&\phantom{-}4t\\
-1,7&=&-5t \quad\text{avec}\: t \in \R\\
-0,7&=& -2t
\end{array}\right.$

Les deux premières équations donnent $t = 0,34$ et la dernière $t = 0,35$. \\
Donc F $\notin \Delta$.
	\end{enumerate}	
\item
	\begin{enumerate}
		\item 
		On a $\vect{\text{AB}}\begin{pmatrix}2\\2\\-1\end{pmatrix}$ et $\vect{\text{AC}}\begin{pmatrix}2\\0\\4\end{pmatrix}$ : ces deux vecteurs ne sont manifestement pas colinéaires, donc les trois points A, B et C définissent un plan.
		\item 
On a $\dfrac13\vect{\text{DE}} \cdot \vect{\text{AB}} = 8 - 10 + 2 = 0$ ;

De même $\dfrac13\vect{\text{DE}} \cdot \vect{\text{AC}} = 8 + 0 - 8 = 0$.

Donc $\vect{\text{DE}}$ est orthogonal à deux vecteurs non colinéaires du plan (ABC) ; il est normal à ce plan.

		\item 
		
D'après la question précédente on sait que :

$M(x~;~y~;~z) \in (\text{ABC}) \iff 4x - 5y - 2z + d = 0$ avec $d \in \R$

Or A$(- 1~;~-1~;~3) \in (\text{ABC})$ donc $4\times (-1) - 5\times (-1)  - 2\times 3  + d = 0$ avec $d \in \R$, soit $- 5 + d = 0 \iff d = 5$.

Conclusion : $M(x~;~y~;~z) \in (\text{ABC}) \iff 4x - 5y - 2z + 5 = 0$.
	\end{enumerate}
\item
	\begin{enumerate}
		\item 
		G$(7~;~-4~;~4) \in \Delta \iff \left\{\begin{array}{r c l}
7 &=& -1 + 4t\\
-4 &=& \phantom{-}6 - 5t \quad\text{avec}\: t \in \R\\
4 &=& \phantom{-}8 - 2t
\end{array}\right.$

Ces trois équations ont pour solution $t = 2$, donc G$(7~;~-4~;~4) \in \Delta$.
		\item 
Le point H appartient à la droite $\Delta$ et au plan (ABC). Donc ses coordonnées $x,\,y,\:z$ vérifient le système :

$\left\{\begin{array}{l c l}
x &=& -1 + 4t\\
y &=& \phantom{-}6 - 5t \quad\text{avec}\: t \in \R\\
z &=& \phantom{-}8 - 2t\\
4x - 5y - 2z + 5 &=& 0
\end{array}\right.$

En replaçant $x,\:y,\: z$ par leurs valeurs en fonction de $t$ dans la dernière équation, on obtient :

$4(- 1 + 4t) - 5(6 - 5t) - 2(8 - 2t) + 5 = 0\iff -4 + 16t - 30 + 25t - 16 + 4t + 5 = 0\iff 45t -45 = 0 \iff 45t = 45 \iff t = 1$.

Les coordonnées de H sont donc $(-1 + 4~;~6 - 5~;~8 - 2)$ soit H(3~;~1~;~6).
		\item 
		
		La distance du point G au plan (ABC) est donc égale à GH.
		
		Or GH$^2 = (3 - 7)^2 + (1 + 4)^2 + (6 - 4)^2 = 16 + 25 + 4 = 45$, d'où \\
		GH $= \sqrt{45} = \sqrt{9 \times 5} = \sqrt{9} \sqrt{5} = 3\sqrt{5}$.
	\end{enumerate}
\item
	\begin{enumerate}
		\item 
D'après la question \textbf{2. a.}, on a $\vect{\text{AB}} \cdot \vect{\text{AC}}  = 2 \times 2 + 2 \times 0 + 4 \times (- 1) = 4 - 4 = 0$ : les vecteurs $\vect{\text{AB}} $ et  $\vect{\text{AC}}$ sont orthogonaux, donc les droites (AB) et (AC) sont perpendiculaires : le triangle ABC est rectangle en A.

		\item 
En prenant comme base le triangle ABC, la hauteur correspondante est GH, donc :

$V = \dfrac13 \times \mathcal{A}(\text{ABC}) \times \text{GH} = \dfrac13 \times \dfrac12 \times \text{AB} \times \text{AC}  \times \text{GH}$.

On a AB$^2 = 2^2 + 2^2 + (- 1)^2 = 4 + 4 + 1 = 9$, donc AB $= 3$ ;

AC$^2 = 2^2 + 0^2 + 4^2 = 4 + 16 = 20$, donc AC $ = \sqrt{20} = \sqrt{4 \times 5} = 2\sqrt{5}$.

Donc $V = \dfrac13 \times \dfrac12 \times 3\times 2\sqrt{5}\times 3\sqrt{5} = 15$.
		

	\end{enumerate}
\end{enumerate}

\end{document}
